% This LaTeX document needs to be compiled with XeLaTeX.
\documentclass[12pt]{article}
\usepackage[margin=0.8in]{geometry}
\usepackage[utf8]{inputenc}
\usepackage{graphicx}
\usepackage[export]{adjustbox}
\graphicspath{ {../images/} }
\usepackage{amsmath}
\usepackage{amsfonts}
\usepackage{amssymb}
\usepackage[version=4]{mhchem}
\usepackage{stmaryrd}
\usepackage{bbold}
%\usepackage[fallback]{xeCJK}
\usepackage{polyglossia}
\usepackage{fontspec}
%\setCJKmainfont{Noto Serif CJK TC}

\usepackage{tikz}
\usepackage{amsthm}
\usepackage{tcolorbox}
\usepackage{enumitem}% http://ctan.org/pkg/enumitem

\tcbuselibrary{theorems}

\definecolor{GrayBG}{HTML}{d9d9d9}
\definecolor{BlueDef}{HTML}{104e8b}
\definecolor{GreenProp}{HTML}{025540}
\definecolor{RedMet}{HTML}{cd5556}
\definecolor{BrownHist}{HTML}{8b5a2b}
\definecolor{BlackSav}{HTML}{000000}

\newtcbtheorem
  []% init options
  {definition}% name
  {Définition}% title
  {%
    colback=GrayBG,
    colframe=BlueDef,
    fonttitle=\bfseries,
  }% options
  {def}% prefix

\newtcbtheorem
  []% init options
  {proposition}% name
  {Proposition}% title
  {%
    colback=GrayBG,
    colframe=GreenProp,
    fonttitle=\bfseries,
  }% options
  {prop}% prefix

\newtcbtheorem
  []% init options
  {method}% name
  {Méthode}% title
  {%
    theorem name,
    colback=GrayBG,
    colframe=RedMet,
    fonttitle=\bfseries,
  }% options
  {met}% prefix
  
\newtcbtheorem
  []% init options
  {history}% name
  {Histoire}% title
  {%
    theorem name,
    colback=GrayBG,
    colframe=BrownHist,
    fonttitle=\bfseries,
  }% options
  {hist}% prefix
  
\newtcbtheorem
  []% init options
  {savoir}% name
  {Savoir-faire du chapitre}% title
  {%
    theorem name,
    colback=GrayBG,
    colframe=BlackSav,
    fonttitle=\bfseries,
  }% options
  {sav}% prefix

\newtheoremstyle{break}
  {}
  {}
  {}
  {}
  {\bfseries}
  {}
  {\newline}
  {}
\theoremstyle{break}
\newtheorem*{example}{Exemple.}
\newtheorem*{remark}{Remarque.}

\newenvironment{demo}[1][\proofname]{%
  \begin{proof}[#1]$ $\newline
}{%
  \end{proof}
}

\setmainlanguage{french}
%\setmainf[tikznode]ont{CMU Serif}

\title{
\vspace{4em}
\centering
\tcbset{colframe=BlackSav,colback=GrayBG}
\tcbox[tikznode]{
Chapitre 4 \\ Arithmétique: Divisibilité dans $\mathbb{Z}$
}
}

\author{}
\date{}


\begin{document}
\maketitle

\tableofcontents

\newpage
\noindent
\begin{minipage}{.7\textwidth}
    \begin{history}{Arithmétique}{}
    L'arithmétique est la science des nombres. Étymologiquement, arithémtique vient en effet du grec arithmos qui signifie nombre. On distingue en général l'arithmétique, basée sur des méthodes relativement élémentaires comme les congruences, et la théorie des nombres faisant intervenir des méthodes plus avancées. Au sein de la théorie des nombres, on distingue également la théorie analytique des nombres (basée sur l'analyse réelle et complexe) et la théorie algébrique des nombres (basée sur l'étude des ensembles de nombres et des structures abstraites). Pour ce qui est de ce dernier domaine, Emmy Nother (1882-1935) a travaillé sur les corps de nombre et les invariants algébriques au début du $\mathrm{xx}^{\mathrm{e}}$ siècle. Elle était considérée par Einstein lui même comme «le génie mathématique créatif le plus considérable produit depuis que les femmes ont eu accès aux études supérieures ». Cela est d'autant plus remarquable qu'en tant que femme juive, elle a dû doublement faire face aux préjugés et aux discriminations de la société allemande de l'époque. En 1933, après l'arrivée d'Hitler au pouvoir, elle se vera d'ailleurs retirer son statut de Professeur et sera exclue de l'Université de Göttingen.
    \end{history}
\end{minipage}%
\begin{minipage}{0.35\textwidth}
    \includegraphics[max width=\textwidth, center]{2024_10_25_a2e0ddf851c541d74f82g-2}\\
\end{minipage}

\section{Relation de divisibilité dans $\mathbb{Z}$}
\subsection{Définition et premières propriétés}
\begin{definition}{}{}
    Soient $a, b \in \mathbb{Z}$. On dit que $a$ divise $b$ lorsqu'il existe $k \in \mathbb{Z}$ tel que $b=a k$. On dit aussi que $a$ est un diviseur de $b$ et que $b$ est un multiple de $a$. On note $a \mid b$.
\end{definition}

\begin{example}
    $3 \mid 15$ car $15=3 \times 5$.
\end{example}

\begin{proposition}{}{}
    Soit $b$ un entier non nul.
    \begin{itemize}[noitemsep,topsep=0pt]
      \item Si $a \mid b$ alors $|a| \leqslant|b|$.
      \item L'entier $b$ admet donc un nombre fini de diviseurs.
    \end{itemize}
\end{proposition}

\begin{proposition}{}{}
    Soient $a, b \in \mathbb{Z}$. Alors,
    
    $$
    a|b \Longleftrightarrow-a| b \Longleftrightarrow a|-b \Longleftrightarrow-a|-b
    $$
\end{proposition}

\begin{demo}
    Montrons que si $a \mid b$ alors $-a \mid b$.\\
    Supposons que $a \mid b$. Par définition, il existe $k \in \mathbb{Z}$ tel que $b=a k$.\\
    Ainsi, $b=(-a) \times(-k)$ et on en déduit donc que $-a \mid b$.\\
    La réciproque et les autres équivalences se montrent d'une manière similaire.
\end{demo}

\subsection{Propriétés de la divisibilité}
\begin{proposition}{Transitivité}{}
    Soient $a, b, c \in \mathbb{Z}$.\\
    Si $a \mid b$ et $b \mid c$, alors $a \mid c$.
\end{proposition}

\begin{demo}
    Supposons que $a \mid b$ et $b \mid c$. Ainsi, il existe $k \in \mathbb{Z}$ tel que $b=a k$. De même, il existe $k^{\prime} \in \mathbb{Z}$ tel que $c=b k^{\prime}$.\\
    Finalement, on en déduit que $c=a k k^{\prime}$ et donc que $a \mid c$.
\end{demo}

\begin{example}
    $3 \mid 9$ et $9 \mid 45$ donc $3 \mid 45$.
\end{example}

\begin{proposition}{}{}
    Soient $a, b, c \in \mathbb{Z}$ tels que $a \mid b$ et $a \mid c$.
    
    \begin{itemize}[noitemsep,topsep=0pt]
      \item Pour tous $m, n \in Z, a \mid(m b+n c)$.
      \item En particulier, $a \mid(b+c)$ et $a \mid(b-c)$.
    \end{itemize}
\end{proposition}

\begin{demo}
    $a \mid b$ et $a \mid c$ donc il existe $k \in \mathbb{Z}$ tel que $b=a k$ et il existe $k^{\prime} \in \mathbb{Z}$ tel que $c=a k^{\prime}$. Ainsi,
    
    $$
    m b+n c=m a k+n a k^{\prime}=a\left(m k+n k^{\prime}\right)
    $$
    \noindent 
    et on en déduit donc que $a \mid(m b+n c)$.
\end{demo}

\begin{example}
    $3 \mid 60$ et $3 \mid 27$ donc $3 \mid 87$
\end{example}

\begin{remark}
Attention! La proprosition «Si $a \mid c$ et si $b \mid c$ alors $a+b \mid c »$ est fausse.\\
Par exemple $2 \mid 6$ et $3 \mid 6$ mais 5 ne divise pas 6 .
\end{remark}

\begin{method}{Déterminer l'ensemble des entiers vérifiant une relation de divisibilité de la forme $(a n+b) \mid(c n+d)$}{}
    \begin{enumerate}[noitemsep,topsep=0pt]
      \item On cherche une combinaison linéaire de $(a n+b)$ et $(c n+d)$ de manière à éliminer l'inconnu $n$.
      \item On en déduit que $a n+b$ est un diviseur de la combinaison linéaire.
      \item On en déduit finalement les valeurs possibles de $a n+b$ puis de $n$.
      \item On vérifie réciproquement que les entiers trouvés conviennent.
    \end{enumerate}
\end{method}

\begin{example}
Déterminer les entiers relatifs $n$ tels que $(2 n+7) \mid(3 n+5)$.\\\\
Solution :\\
Soit $n \in \mathbb{Z}$ tel que $(2 n+7) \mid(3 n+5)$.\\
On a $3 \times(2 n+7)-2 \times(3 n+5)=11$.\\
Ainsi, comme $2 n+7 \mid 2 n+7$ et $2 n+7 \mid 3 n+5$, on en déduit que $2 n+7 \mid 11$.\\
Par conséquent : $2 n+7=11$ ou $2 n+7=1$ ou $2 n+7=-11$ ou $2 n+7=-1$.\\
Finalement, on obtient donc $n=2$ ou $n=-3$ ou $n=-9$ ou $n=5$.\\
Réciproquement, on vérifie que $2,-3,-9$ et 5 sont bien des solutions.\\
En conclusion,

$$
\mathcal{S}=\{2 ;-3 ;-9 ; 5\}
$$
\end{example}

\subsection{Division euclidienne}
\begin{proposition}{Division euclidienne dans $\mathbb{N}$}{}
    Soient $a \in \mathbb{N}$ et $b \in \mathbb{N} \backslash\{0\}$.\\
    Il existe un unique couple $(q, r) \in \mathbb{N} \times \mathbb{N}$ tels que $a=b q+r$ et $0 \leqslant r<b$. L'entier $q$ est appelé le quotient de la division de $a$ par $b$ et $r$ le reste.
\end{proposition}

\begin{remark}
    Pour démontrer cette proposition, on utilisera le fait que tout ensemble $A$ inclus dans $\mathbb{N}$ et non vide admet un minimum (c'est-à-dire qu'il existe $m \in A$ tel que pour tout $n \in A, m \leqslant n$ ).
\end{remark}
\begin{demo}
    Soient $a \in \mathbb{N}$ et $b \in \mathbb{N} \backslash\{0\}$.
    
    \begin{itemize}[noitemsep,topsep=0pt]
        \item Montrons tout d'abord l'existence des entiers $q$ et $r$.
    
        On définit l'ensemble $A=\{a-b n, n \in \mathbb{N}\} \bigcap \mathbb{N}$. Alors $A$ est inclus dans $\mathbb{N}$ et est non vide (pour $n=0$, on voit que $a \in A$ ).\\
        Ainsi, on en déduit que $A$ admet un élément minimum que l'on note $r$.\\
        Comme $r \in A$, par définition de $A$, il existe $q \in \mathbb{N}$ tel que $a-b q=r$.\\
        Finalement, on a prouvé qu'il existe des entiers $a$ et $r$ tels que $a=b q+r$. De plus, suppsoons par l'absurde que $r \geqslant b$. En posant $q^{\prime}=q+1$, on aurait $a=b q^{\prime}+r-b$ et par conséquent, $r-b \in A$. Cela est absurde car $r$ est l'élément minimum de $A$.
    
        \item Montrons l'unicité du couple $(q, r)$.
    
        On suppose qu'il existe deux couples $\left(q_{1}, r_{1}\right)$ et $\left(q_{2}, r_{2}\right)$ vérifiant les conditions de la propriété.\\
        Ainsi, $a=b q_{1}+r_{1}=b q_{2}+r_{2}$.\\
        On en déduit que $b\left(q_{1}-q_{2}\right)=r_{2}-r_{1}$ et donc que $b \mid r_{2}-r_{1}$. Par ailleurs, on sait que $0 \leqslant r_{1}<b$ et $0 \leqslant r_{2}<b$. Par conséquent, $-b<r_{2}-r_{1}<b$. Comme on a montré par ailleurs que $b \mid r_{2}-r_{1}$, on en déduit que $r_{2}-r_{1}=0$, c'est-à-dire que $r_{2}=r_{1}$. Par suite, on obtient aussi que $q_{2}=q_{1}$ ce qui prouve l'unicté.
    \end{itemize}
\end{demo}

\begin{proposition}{Division euclidienne dans $\mathbb{Z}$ (admise)}{}
    Soient $a, \in \mathbb{Z}$ et $b, \in \mathbb{Z} \backslash\{0\}$.\\
    Il existe un unique couple $(q, r) \in \mathbb{Z} \times \mathbb{N}$ tels que $a=b q+r$ et $0 \leqslant r<|b|$.
\end{proposition}

\begin{example}
    \begin{itemize}[noitemsep,topsep=0pt]
        \item[]
        \item La division euclidienne de 35 par 6 est $35=6 \times 4+3$.
        \item La division euclidienne de -15 par 4 est $-15=4 \times(-4)+1$.
    \end{itemize}
\end{example}

\noindent
\begin{minipage}{.7\textwidth}
    \begin{history}{Division euclidienne}{}
        On parle de «division euclidienne» en référence au mathématicien grec Euclide (300 av. J.C.). La division était néanmoins connue bien avant lui et utilisée par exemple par les égyptiens. Le nom «euclidien» s'explique plutôt par le fait qu'Euclide a utilisé un algorithme basé sur la division euclidienne (voir chapitre suivant). Il le présente dans un ouvrage célèbre intitulé Les Éléments. Avec plus de mille éditions différentes depuis l'antiquité, il s'agit, juste après la Bible, de l'ouvrage le plus imprimé dans l'histoire de l'humanité.
    \end{history}
\end{minipage}%
\hspace{1em}
\begin{minipage}{0.35\textwidth}
    \begin{tcolorbox}[colback=white,colframe=BrownHist]
        \centering
        \includegraphics[max width=\textwidth, center]{2024_10_25_a2e0ddf851c541d74f82g-5}
        Euclide
    \end{tcolorbox}
\end{minipage}

\section{Congruences}
\subsection{Définition}
\begin{definition}{}{}
    Soit $m \in \mathbb{N}^{*}$ et $a, b \in \mathbb{Z}$.\\
    On dit que $a$ et $b$ sont congrus modulo $\mathbf{m}$ s'ils ont le même reste dans la division euclidienne par $m$.\\
    On note $a \equiv b[m]$ ou $a \equiv b \bmod m$.
\end{definition}

\begin{example}
    Le reste de la division de 16 par 5 est 1 . De même, le reste de la division de 31 par 5 est 1. Ainsi, $16 \equiv 31[5]$.
\end{example}

\begin{proposition}{}{}
    Soit $m \in \mathbb{N}^{*}$ et $a, b \in \mathbb{Z}$.
    $$
    a \equiv b[m] \Longleftrightarrow m \mid(b-a)
    $$
\end{proposition}

\begin{demo}
    Soit $m \in \mathbb{N}^{*}$ et $a, b \in \mathbb{Z}$.
    
    \begin{itemize}
        \item Supposons que $a \equiv b[m]$. Cela signifie que $a$ et $b$ ont le même reste dans la division euclidienne par $m$. Autrement dit, il existe $p, q, r \in \mathbb{Z}$ tels que $a=m p+r$ et $b=m q+r$ avec $0 \leqslant r<|m|$.\\
        Ainsi, $b-a=(m q+r)-(m p+r)=m \times(q-p)$.\\
        Cela montre bien que $m$ divise $b-a$.
        \item Réciproquement, supposons que $m \mid(b-a)$.\\
        On effectue alors les divisions euclidiennes de $a$ et $b$ par $m$.\\
        D'une part, il existe $q, r \in \mathbb{Z}$ tels que $a=m q+r$ avec $0 \leqslant r<|m|$.\\
        D'autre part, il existe $q^{\prime}, r^{\prime} \in Z$ tels que $b=m q^{\prime}+r^{\prime}$ avec $0 \leqslant r^{\prime}<|m|$.\\
        L'objectif est de démontrer alors que $r=r^{\prime}$.\\
        En fait, $b-a=\left(m q^{\prime}+r^{\prime}\right)-(m q+r)=m\left(q^{\prime}-q\right)+\left(r^{\prime}-r\right)$.\\
        Or, les inégalités vérifiées par $r$ et $r^{\prime}$ impliquent nécessairement que $\left(r^{\prime}-r\right) \leqslant|m|$.\\
        Ainsi, si $r^{\prime}-r \geqslant 0$, l'égalité 1 correspond à la division euclidienne de $b-a$ par $m$.\\
        Mais comme on sait que $m \mid(b-a)$, on en déduit que le reste $r^{\prime}-r$ est nul, c'est-à-dire, $r^{\prime}=r$.\\
        Autrement dit, $a$ et $b$ ont même reste dans la division par $m$ et sont donc congrus modulo $m$.\\
        Dans le cas où $r^{\prime}-r<0$, il suffit de raisonner de même avec $a-b$ au lieu de $b-a$.
    \end{itemize}
\end{demo}

\noindent
\begin{minipage}{.7\textwidth}
    \begin{history}{Congruences}{}
        Carl Friedrich Gauß (1777-1855) est originaire d'une famille pauvre de la principauté de Brunswick. Dès l'école primaire, l'instituteur et son assistant décèlent ses talents et lui transmettent leur passion pour les mathématiques. Gauß publie ses premiers résultats à 19 ans et à 24 ans, il introduit la notion de congruences dans un ouvrage célèbre intitulé Disquisitiones Arithmeticae. Dès 28 ans, il dirigea l'observatoire astronomique de Göttingen. Alors qu'il n'aimait pas vraiment enseigner, cela lui permettait justement de se consacrer à ses recherches mathématiques.
    \end{history}
\end{minipage}%
\hspace{1em}
\begin{minipage}{0.35\textwidth}
    \begin{tcolorbox}[colback=white,colframe=BrownHist]
        \centering
        \includegraphics[max width=\textwidth, center]{2024_10_25_a2e0ddf851c541d74f82g-6}
        Carl Friedrich Gauß
    \end{tcolorbox}
\end{minipage}

\subsection{Congruences et opérations}
\begin{proposition}{Transitivité}{}
    Soit $m \in \mathbb{N}^{*}$ et $a, b, c \in \mathbb{Z}$.
    \begin{center}
        Si $a \equiv b[m]$ et $b \equiv c[m]$ alors $a \equiv c[m]$
    \end{center}
\end{proposition}

\begin{demo}
    Évident en revenant à la définition d'une congruence.
\end{demo}

\begin{proposition}{}{}
    Soit $m \in \mathbb{N}^{*}$ et $a, b, c, d \in \mathbb{Z}$.
    \begin{itemize}[noitemsep,topsep=0pt]
        \item \textbf{Compatibilité avec l'addition :}\\
        Si $a \equiv b[m]$ et $c \equiv d[m]$, alors $a+c \equiv b+d[m]$
        \item \textbf{Compatibilité avec la multiplication :}\\
        Si $a \equiv b[m]$ et $c \equiv d[m]$, alors $a \times c \equiv b \times d[m]$
        \item \textbf{Compatibilité avec les puissances :}\\
        Si $a \equiv b[m]$, alors pour tout $p \in \mathbb{N}^{*}, a^{p} \equiv b^{p}[m]$
    \end{itemize}
\end{proposition}

\begin{demo}
    Soient $a \equiv b[m]$ et $c \equiv d[m]$.\\
    Cela signifie qu'il existe $k \in \mathbb{Z}$ tel que $a-b=m k$ et il existe $k^{\prime} \in \mathbb{Z}$ tel que $c-d=m k^{\prime}$.\\
    Autrement dit, $a=b+m k$ et $c=d+m k^{\prime}$.\\
    En sommant les deux égalités, on obtient : $a+c=b+d+m\left(k+k^{\prime}\right)$, ce qui signifie exactement que $a+c \equiv b+d[m]$.\\
    De la même manière, en multipliant les deux égalités, $a \times c=(b+m k) \times\left(d+m k^{\prime}\right)=b \times d+$ $m\left(d k+b k^{\prime}\right)$.\\
    On en déduit que $a \equiv b[m]$ et $c \equiv d[m]$, alors $a \times c \equiv b \times d[m]$.\\
    Finalement, la dernière propriété s'obtient par récurrence comme conséquence de la seconde.
\end{demo}

\begin{method}{Montrer une relation de divisibilité}{}
    \begin{enumerate}[noitemsep,topsep=0pt]
      \item On travaille modulo le diviseur (3 dans l'exemple ci-dessous).
      \item On distingue les cas grâce à un tableau en montrant utilisant le fait qu'être divisible par $n$ est équivalent au fait d'être égal à 0 modulo $n$.
    \end{enumerate}
\end{method}

\begin{example}
Montrer que pour tout entier $n \in \mathbb{Z}, 2 n(n+1)(n+5)$ est divisible par 3.\\\\
Solution :\\
On distingue les cas modulo 3 :
\end{example}

\begin{center}
\begin{tabular}{|c|c|c|c|}
\hline
$n$ & 0 & 1 & 2 \\
\hline
$2 n$ & 0 & 2 & 1 \\
\hline
$n+1$ & 1 & 2 & 0 \\
\hline
$n+5$ & 2 & 0 & 1 \\
\hline
$2 n(n+1)(n+5)$ & 0 & 0 & 0 \\
\hline
\end{tabular}
\end{center}
\noindent
Ainsi, dans tous les cas, $2 n(n+1)(n+5) \equiv 0[3]$ ce qui signifie exactement que $2 n(n+1)(n+5)$ est divisible par 3.

\subsection{Inverse modulo $m$}
\begin{definition}{}{}
    Soient $m \in \mathbb{Z}^{*}$ et $a \in \mathbb{Z}$.\\
    On dit que $a$ est inversible modulo $m$ lorsqu'il existe un entier $b$ tel que $a \times b \equiv 1[m]$. De plus, l'entier $b$ est appelé \textbf{inverse de $\boldsymbol{a}$ modulo $\boldsymbol{m}$}.
\end{definition}

\begin{example}
    8 est inversible modulo 3 car $8 \times 2 \equiv 1[3]$. Son inverse est 2 modulo 3 .
\end{example}

\begin{proposition}{Unicité de l'inverse}{}
    Soient $m \in \mathbb{Z}^{*}$ et $a \in \mathbb{Z}$.\\
    Si $a$ est inversible modulo $m$ alors l'inverse est unique modulo $m$.
\end{proposition}

\begin{demo}
    Supposons qu'il existe deux inverse de $a$, notés $b_{1}$ et $b_{2}$. On va montrer que $b_{1} \equiv b_{2}[m]$.\\
    En fait, $a \times b_{1} \equiv 1[m]$ (car $b_{1}$ est un inverse).\\
    Ainsi, en multipliant par $b_{2}$, il vient :\\
    $b_{2} \times a \times b_{1} \equiv b_{2} \times 1[m]$\\
    donc $1 \times b_{1} \equiv b_{2}[m]$\\
    donc $b_{1} \equiv b_{2}[m]$\\
    Cela prouve donc l'unicité, modulo $m$ de l'inverse.
\end{demo}

\begin{method}{Déterminer si un entier est inversible ou non}{}
    On établit un tableau de congruences afin de distinguer tous les cas.
\end{method}

\begin{example}
\begin{enumerate}[noitemsep,topsep=0pt]
    \item[]
    \item Montrer que 7 est inversible modulo 9.
    \item Montrer que 4 n'est pas inversible modulo 6 .
\end{enumerate}
\noindent
~\\Solution :
\begin{enumerate}[noitemsep,topsep=0pt]
    \item
    \begin{tabular}{|c|c|c|c|c|c|c|c|c|c|}
    \hline
    $n \equiv \ldots[9]$ & 0 & 1 & 2 & 3 & 4 & 5 & 6 & 7 & 8 \\
    \hline
    $7 n \equiv \ldots[9]$ & 0 & 7 & 5 & 3 & 1 & $\ldots$ & $\ldots$ & $\ldots$ & $\ldots$ \\
    \hline
    \end{tabular}
    
    Ainsi, $7 \times 4 \equiv 1[9]$ donc 7 est inversible modulo 9.\\
    \item
    \begin{tabular}{|c|c|c|c|c|c|c|}
    \hline
    $n \equiv \ldots[6]$ & 0 & 1 & 2 & 3 & 4 & 5 \\
    \hline
    $4 n \equiv \ldots[6]$ & 0 & 4 & 2 & 0 & 4 & 2 \\
    \hline
    \end{tabular}
    
    Ainsi, pour tout entier $n, 4 n \not \equiv 1[6]$, ce qui signifie que 4 n'est pas inversible modulo 6 .
\end{enumerate}
\end{example}

\subsection{Application à l'étude des règles de divisibilité}\noindent
Dans cette partie, on se servira du fait que tout nombre entier $N$ peut s'écrire sous la forme $N=\sum_{k=0}^{n} a_{k} 10^{k}$, où pour tout $k \in\{1, \ldots, n\}, 0 \leqslant a_{k} \leqslant 9$ et $a_{n} \neq 0$.

\begin{proposition}{}{}
    Un nombre entier $N$ est divisible par 10 si, et seulement si, son chiffre des unités est 0 .
\end{proposition}

\begin{demo}
    En écrivant $N=\sum_{k=0}^{n} a_{k} 10^{k}$, cela revient à prouver que $N$ est divisible par 10 si, et seulement si, $a_{0}=0$.\\
    En fait, comme pour tout $k \geqslant 1,10^{k} \equiv 0[10]$, on a : $N \equiv a_{0}[10]$, d'où le résultat.
\end{demo}

\begin{proposition}{}{}
    Un nombre entier $N=\sum_{k=0}^{n} a_{k} 10^{k}$ est divisible par 3 si, et seulement si, $a_{0}+a_{1}+\ldots+a_{n}$ est divisible par 3 .
\end{proposition}

\begin{demo}
    Comme $10 \equiv 1[3]$, on en déduit que $10^{k} \equiv 1[3]$ (d'après la compatibilité des congruences avec les puissances).\\
    Finalement, par compatibilité des congruences avec la multiplication, on en déduit que :
    
    $$
    N \equiv \sum_{k=0}^{n} a_{k} 10^{k} \equiv \sum_{k=0}^{n} a_{k} \times 1 \equiv \sum_{k=0}^{n} a_{k}[3]
    $$
\end{demo}

\begin{savoir}{}{}
    \begin{itemize}[noitemsep,topsep=0pt]
      \item Déterminer l'ensemble des diviseurs d'un entier.
      \item Déterminer l'inverse (lorsqu'il existe) d'un nombre entier modulo $n$.
      \item Déterminer l'ensemble des entiers vérifiant une relation de divisibilité de la forme $(a n+b) \mid(c n+d)$
      \item Utiliser les congruences pour établir une relation de divisibilité.
      \item Utiliser un raisonnement par disjonction de cas (pour établir une divisibilité par exemple).
    \end{itemize}
\end{savoir}

\end{document}